\section{Admissibility from diagonal local parameters}\label{sec:diagonal}

Proposition~\ref{prop:family} gives an admissible compact object for
\(G^b\), where \(G=\GL_2(\Q_2)\) and \(b=[F:\Q]\). We need an
admissible object for one copy of \(G\), still carrying the whole trace
family \(D\). The equality of the \(b\) local Galois parameters is the
additional information that makes this possible.

Retain the local block generator \(P\), its endomorphism algebra
\(\mathcal E=\End(P)\), and its fixed-determinant pseudodeformation ring
\(R\) from Proposition~\ref{prop:family}. Thus the central
character is \(\psi=\chi\eps\), the Galois determinant in block
normalization is \(\psi\eps\), and the compact category uses the dual
central action. Write
\[
 \mathcal E^{(b)}=\widehat\bigotimes_{i=1}^b\mathcal E,
 \qquad P^{(b)}=\widehat\bigotimes_{i=1}^bP.
\]
All completed tensor products in this section are in the category of
pseudocompact modules. Coinvariants mean quotients by the closed
submodule generated by the indicated relations.

\begin{proposition}[One-factor admissibility]\label{prop:diagonal}
Let \(D\) be a complete noetherian local \(\OO\)-domain in which
\(\varpi\ne0\), and suppose that \(R\to D\) is a finite continuous
local map. Let \(L\ne0\)
be a finite torsion-free \(D\)-module with a commuting continuous right
\(\mathcal E^{(b)}\)-action. Suppose that every copy of \(R\) acts
through the same map \(R\to D\), and that
\[
 M=L\widehat\otimes_{\mathcal E^{(b)}}P^{(b)}
\]
is admissible in the compact category. Then, using the first
\(\mathcal E\)-action,
\[
 V=L\widehat\otimes_{\mathcal E}P
\]
is nonzero, \(\OO\)-torsion-free, and finite over
\(\mathscr A=\OO[[K]]\), for every sufficiently small uniform open
subgroup \(K\subset\SL_2(\Z_2)\). It has a continuous commuting
\(D\)-action.
\end{proposition}

The construction takes place on the algebra side of the block equivalence:
\[
\begin{CD}
 L_{\mathcal E^{(b)}} @>{-\widehat\otimes_{\mathcal E^{(b)}}P^{(b)}}>> M\\
 @V{\text{forget factors }2,\ldots,b}VV @.\\
 L_{\mathcal E} @>{-\widehat\otimes_{\mathcal E}P}>> V.
\end{CD}
\]
The subscripts indicate which algebra acts on \(L\). The lower row
reconstructs the one-factor object from that algebra module.

Here admissibility in the compact category means finite generation over
the corresponding completed compact group algebra. Although \(K\) is
not open in \(G\), its product with the scalar units is open in
\(\GL_2(\Z_2)\); the fixed central character therefore makes this the
usual admissibility condition. We prove the proposition by comparing
one-factor and product coinvariants. The essential local assertion will
be that, along a nonconstant characteristic-two curve, all simple block
modules give the same answer to the question whether their
\(K\)-coinvariants vanish.

\subsection{The local block input}

We record the precise consequences of Pa\v{s}k\=unas--Tung that enter the
argument. The residue field has been enlarged so that the block and its
simple objects are absolutely split. Let \(e\in\mathcal E\) be the sum
of the idempotents corresponding to those simple objects on which
\(\SL_2(\Q_2)\) does not act trivially.

\begin{theorem}[Pa\v{s}k\=unas--Tung]\label{input:local-blocks}
The following statements hold for every such block at \(p=2\),
including the block of a scalar residual Galois representation.
\begin{enumerate}
\item The functor \(\Hom(P,-)\) is an exact equivalence from the compact
block to right pseudocompact \(\mathcal E\)-modules, with inverse
\(-\widehat\otimes_{\mathcal E}P\). The algebra \(\mathcal E\) is finite
over \(R\).
\item Let \(\mathrm{CH}_R\) be the universal Cayley--Hamilton algebra
for the local pseudorepresentation, and let
\(\mathrm{CH}_R^{\mathrm{tf}}\) be its maximal
\(\OO\)-torsion-free quotient. There is an \(R\)-compatible
identification
\[
 \mathrm{CH}_R^{\mathrm{tf}}\simeq
 \begin{cases}
  (e\mathcal E e)^{\mathrm{op}},&\text{in a nonsupersingular block},\\
  M_2((e\mathcal E e)^{\mathrm{op}}),&\text{in a supersingular block}.
 \end{cases}
\]
The corner equivalence is induced by the Montr\'eal functor on the
quotient by the objects with trivial \(\SL_2(\Q_2)\)-action.
\item An \(\mathcal E\)-module finite and torsion-free over \(\OO\) gives,
after duality and inversion of \(2\), an admissible Banach representation
of finite length. At an absolutely irreducible Galois coefficient
parameter there is a unique absolutely irreducible Banach object. At a
sum of characters \(\gamma_1+\gamma_2\), the simple objects have the
following description. If the characters coincide, the unique simple is
\(\Ind_B^G(1\otimes\eps^{-1})^{\cts}\otimes\gamma_1\circ\det\).
If their sum is \(\gamma+\gamma\eps\), the simples are
\[
 \{1,\widehat{\mathrm{St}},
       \Ind_B^G(\eps\otimes\eps^{-1})^{\cts}\}
       \otimes\gamma\circ\det,
\]
where \(\widehat{\mathrm{St}}\) is continuous Steinberg. If
\(\gamma_1\gamma_2^{-1}\notin\{1,\eps,\eps^{-1}\}\), they are exactly
\[
 \Ind_B^G(\gamma_1\otimes\gamma_2\eps^{-1})^{\cts},
 \qquad
 \Ind_B^G(\gamma_2\otimes\gamma_1\eps^{-1})^{\cts}.
\]
Here \(B\) is the upper Borel, induction is unnormalized, and Galois
characters are viewed as characters of \(\Q_2^\times\) by reciprocity.
\end{enumerate}
\end{theorem}

The first assertion is \cite[Section~4.1 and Theorems~1.2--1.3]
{PaskunasTung2021}; the second is
\cite[Propositions~4.10 and~4.18, Theorem~6.13]{PaskunasTung2021};
the Banach assertions are \cite[Section~4.5, Section~6.2,
Equation~(20), Proposition~6.9, Corollary~6.10, and
Proposition~6.11]{PaskunasTung2021}.
At \(2\), the center comparison gives an isomorphism after inverting
\(2\) and permits an integral cokernel killed by \(2\)
\cite[Theorem~1.4]{PaskunasTung2021}. We use the finite map from \(R\)
and the stated integral Cayley--Hamilton comparison, without identifying
\(R\) with the integral center.

The compact equivalence commutes with coefficient quotients, finite
presentations, and inverse limits. Consequently it preserves injections
and coefficient torsion-freeness, and commutes with the completed
coefficient changes below. The corresponding product equivalence is
explained in \cite[Proposition~4.2]{FM2}.

\subsection{Coinvariants and the noncharacter corner}

Fix a sufficiently small uniform \(K\subset\SL_2(\Z_2)\), and put
\[
 C_K=P/(\varpi P+(K-1)P).
\]
This is a left \(\mathcal E\)-module. Its support over \(R\) will control
the admissibility test.

\begin{lemma}\label{lem:coinvariant-finite}
The module \(C_K\) is finite over \(R\). For a finite pseudocompact
right \(\mathcal E\)-module \(W\),
\[
 (W\widehat\otimes_{\mathcal E}P)/(\varpi,K-1)
       =W\otimes_{\mathcal E}C_K.
\]
The analogous formula holds for product coinvariants. If a commutative
complete coefficient algebra $A$, finite over $R$ (over
$R^{\widehat\otimes b}$ in the product case), acts on $W$ compatibly
and commutes with its algebra action, these formulas commute with
base change $A\to H$ to a field.
\end{lemma}
\begin{proof}
The object \(P/\mathfrak m_RP\) corresponds under the compact
equivalence to \(\mathcal E/\mathfrak m_R\mathcal E\), so has finite
length. Every smooth simple in the block is admissible. Its compact
dual therefore has finite-dimensional \(K\)-coinvariants, with the
central character understood. Right exactness of coinvariants now
makes \(C_K/\mathfrak m_RC_K\) finite-dimensional. The
\(\mathfrak m_R\)-action is topologically nilpotent: on the algebra
side this follows from the finite \(R\)-module \(\mathcal E\), and the
equivalence transports the inverse limit of its finite-length
quotients. Compact Nakayama proves that \(C_K\) is finite over \(R\).

Completed tensor products commute with the cokernels defining
coinvariants. After these cokernels, all modules involved are finite
over their complete noetherian coefficient rings, so ordinary and
completed tensor products agree. For the product formula, apply this
argument in each factor. Finite presentations over \(R\), or over its
completed tensor powers, then show that the formula commutes with the
asserted field base changes.
\end{proof}

A characteristic-two curve in \(\Spec R\) means the spectrum of a
one-dimensional complete local domain quotient of \(R/\varpi\). Its
generic point is nonconstant if its map to \(\Spec R\) is not the
closed residual point.

\begin{lemma}\label{lem:full-corner}
At the generic point of a nonconstant characteristic-two curve, the
idempotent \(e\) is full: if \(H\) is a finite extension of the curve's
fraction field and \(\mathcal E_H=\mathcal E\otimes_RH\), then
\(\mathcal E_He\mathcal E_H=\mathcal E_H\).
\end{lemma}
\begin{proof}
We show that \(\mathcal E/\mathcal E e\mathcal E\) is finite over
\(\OO\). Under the compact equivalence its modules have only the
\(\SL_2(\Q_2)\)-trivial simple constituents. This is precisely the
subcategory of objects with trivial \(\SL_2(\Q_2)\)-action: that
subcategory is thick, and the assertion passes to compact limits
\cite[Section~4.2]{PaskunasTung2021}.

Such an action factors through \(\det:G\to\Q_2^\times\). On squares,
the character is prescribed by the central action. Since
\(\Q_2^\times/(\Q_2^\times)^2\) is finite, the corresponding completed
group algebra with these scalar relations is finite over \(\OO\).
This remains true if the scalar relations force coefficient torsion.
The regular module of \(\mathcal E/\mathcal E e\mathcal E\) corresponds
to an object of this subcategory with finite cosocle. Compact Nakayama
makes it finite over that finite \(\OO\)-algebra; its endomorphism
algebra, and hence \(\mathcal E/\mathcal E e\mathcal E\), is finite over
\(\OO\). Its special-fiber support over \(R\) is therefore contained in the
closed point. It vanishes at the curve's generic point, as required.
\end{proof}

\subsection{The characteristic-two curve test}

Let $H$ be a finite extension of the fraction field of a nonconstant
characteristic-two curve, chosen to split the parameter and the algebra
$\mathcal E_H=\mathcal E\otimes_R H$. We will show that a nonzero
$C_{K,H}=C_K\otimes_R H$ survives tensoring with every simple right
$\mathcal E_H$-module. The main case is a parameter
with two distinct characters. Its two principal-series families have
the same coinvariant vanishing: exchanging their inducing characters
inverts their values on determinant-one stabilizers. For a unit $a$,
the elements $a-1$ and $a^{-1}-1$ generate the same ideal.
To apply this calculation to all block simples, we construct the
families integrally and track their two character labels through
specialization.

\begin{lemma}\label{lem:principal-families}
Suppose a nonconstant characteristic-two curve has generic local
parameter \(\gamma_1+\gamma_2\), with \(\gamma_1\ne\gamma_2\). After a
finite extension, let \(J=\kk'[[t]]\) be its normalization, finite
over the curve ring, and
\(H=\Frac(J)\). The compact dual principal-series families over \(J\)
with inducing orders
\[
 (\gamma_1,\gamma_2\eps^{-1}),\qquad
 (\gamma_2,\gamma_1\eps^{-1})
\]
correspond to finite free \(J\)-modules \(W_1,W_2\) for
\(\mathcal E\). Their central \(R\)-action is the given curve parameter.
Their generic fibers are the two distinct simple right
\(\mathcal E_H\)-modules, and
\[
 (W_1\otimes_JH)\otimes_{\mathcal E_H}(C_K\otimes_RH)\ne0
 \quad\Longleftrightarrow\quad
 (W_2\otimes_JH)\otimes_{\mathcal E_H}(C_K\otimes_RH)\ne0.
\]
\end{lemma}
\begin{proof}
The characters take values in \(J^\times\): their values are integral
by the characteristic polynomial, and their product is a unit. They
are continuous. Indeed, choose an element on which the two characters
differ; traces against that element recover their values continuously
in \(H\). Their reductions are the residual constituents.

We first construct the families in mixed characteristic. Enlarge the
integer coefficients to \(\OO'\), and set
\(J_0=\OO'[[X,Y]]\). The pro-\(2\) completion of
\(\Q_2^\times\) has two free generators, represented by \(2\) and
\(5\), and torsion generated by \(-1\). Let \(\Gamma_1\) be the
universal lift of the first residual character on the two free
generators, fixing either lift of its sign on \(-1\), and put
\(\Gamma_2=\psi\eps\Gamma_1^{-1}\). There is a continuous map
\(J_0\to J\) specializing these characters to \(\gamma_1,\gamma_2\).

For each order, form smooth induction with coefficients in the
discrete dual \(J_0^\vee\), and take its compact dual \(N_i\).
This defines the family even when the inducing characters become
smooth only modulo powers of the coefficient maximal ideal.
Local sections on \(B\backslash G=\mathbf P^1(\Q_2)\) identify its
compact model with a pro-free \(J_0\)-module. In particular the model
is flat and commutes with coefficient changes. The residual smooth
induction has finite length in the specified block. If the residual
characters differ it is irreducible; if they coincide it has the
character and Steinberg constituents. In either case it has precisely
one noncharacter constituent, with multiplicity one
\cite[Section~4.1]{PaskunasTung2021}. Filtering the finite coefficient
quotients shows that \(N_i\) belongs to the compact block.

Put \(\widetilde W_i=\Hom(P,N_i)\). Compact Nakayama and residual
finite length show that this is a finite \(J_0\)-module. Exactness of
the equivalence makes it flat: for every finitely generated ideal
\(I\subset J_0\), apply the equivalence to the injective map
\(I\otimes_{J_0}N_i\to N_i\), using finite presentations to commute
these tensors with the equivalence. Thus \(\widetilde W_i\) is finite
free. The same argument proves compatibility with coefficient changes.
The residual multiplicity calculation gives
\[
                    \rank_{J_0}(\widetilde W_i e)=1.
\]

At characteristic-zero coefficient points outside the proper loci
\(\Gamma_1\Gamma_2^{-1}\in\{1,\eps,\eps^{-1}\}\), Schikhof duality
identifies these objects with the two continuous principal series in
Theorem~\ref{input:local-blocks}. Those points are Zariski dense in
\(\Spec J_0\). Hence their central \(R\)-action is everywhere the
pseudorepresentation \(\Gamma_1+\Gamma_2\): equality may be tested on
matrices acting on the finite free modules.

We also need to keep the two character labels distinct upon
specialization. The residual block here is nonsupersingular, so
Theorem~\ref{input:local-blocks} identifies the opposite corner with
\(\mathrm{CH}_R^{\mathrm{tf}}\), without a matrix factor.
Its Galois action on each free rank-one module
\(\widetilde W_i e\) gives a character. At the dense coefficient
points the two principal series are distinct simple modules, their
nonzero corner modules are distinct, and their corner labels are
therefore \(\Gamma_1\) and \(\Gamma_2\). More explicitly, if
\(\lambda_i\) is the corner character, set
\[
 I_{ij}=(\lambda_i(g)-\Gamma_j(g):g\in G_{\Q_2})\subset J_0.
\]
The dense coefficient points lie in
\(V(I_{i1})\cup V(I_{i2})=V(I_{i1}I_{i2})\), so
\(I_{i1}I_{i2}=0\). Since \(J_0\) is a domain, one of the two ideals
vanishes. The two labels are opposite because the generic
specializations are distinct. Thus the labels on the entire family
are \(\Gamma_1,\Gamma_2\), up to one fixed permutation, and these
identities persist under \(J_0\to J\).

Set \(W_i=\widetilde W_i\otimes_{J_0}J\). By
Lemma~\ref{lem:full-corner}, \(e\) is full over \(H\); its rank-one
corner modules therefore make \(W_i\otimes_JH\) simple and distinct.
The specialized torsion-free Cayley--Hamilton algebra is a quotient
of the Cayley--Hamilton algebra of the specialized parameter. The
field structure theorem for a split multiplicity-free determinant
allows at most the two character simples
\cite[Theorems~2.12 and~2.22]{Chenevier2014}. Thus our two modules
account for all simples, without assuming that forming the torsion-free quotient
commutes with reduction modulo \(2\).

Finally calculate \(K\)-coinvariants directly in the compact
induction model. There are finitely many double cosets
\(B\backslash G/K\). For a representative \(g\), its contribution
is \(J\) modulo the ideal generated by the inducing character minus
one on \(B\cap gKg^{-1}\). This follows equally by dualizing the
invariant functions with values in \(J^\vee\). Every stabilizer has
determinant one. Since \(\eps=1\) over \(J\), the two inducing
characters on that stabilizer are inverse. They generate the same
ideal, because \(a^{-1}-1=-a^{-1}(a-1)\) for a unit \(a\).
The two coinvariant modules consequently vanish simultaneously after
tensoring with \(H\). Lemma~\ref{lem:coinvariant-finite} identifies
these modules with the two displayed tensor products.
\end{proof}


\begin{lemma}[Uniform survival on curves]\label{lem:curve-support}
Let \(H\) be a finite extension of the fraction field of a nonconstant
characteristic-two curve in \(\Spec R\), large enough to split the
parameter and the finite algebra \(\mathcal E_H\). If
\(C_{K,H}=C_K\otimes_RH\ne0\), then
\[
                  S\otimes_{\mathcal E_H}C_{K,H}\ne0
\]
for every simple right \(\mathcal E_H\)-module \(S\).
\end{lemma}
\begin{proof}
The corner is full by Lemma~\ref{lem:full-corner}. The
Cayley--Hamilton comparison and its field structure therefore give
at most the simples labelled by the absolute constituents of the
parameter. Repeated characters cannot occur on this curve. Indeed
the fixed determinant has constant finite values in characteristic
two; if it equals \(\gamma^2\), injectivity of squaring in a field
makes \(\gamma\), and hence the whole parameter, constant.

If the parameter is absolutely irreducible, the split algebra has only
one simple module. A nonzero finite left module has a nonzero
semisimple head, so tensoring it with that simple right module is
nonzero. If the parameter has two distinct characters,
Lemma~\ref{lem:principal-families} constructs both simple
modules and proves that their tensor tests have identical vanishing.
At least one test is nonzero: otherwise the semisimple head of \(C_{K,H}\) would be
zero. Both tests are therefore nonzero.
\end{proof}

\subsection{From product coinvariants to one-factor admissibility}

\begin{proof}[Proof of Proposition~\ref{prop:diagonal}]
The module \(L\) is finite over \(R\), hence over the first
\(\mathcal E\). The exact compact equivalence makes \(V\) nonzero and
\(\OO\)-torsion-free. It also transports the continuous commuting
\(D\)-action. By Lemma~\ref{lem:coinvariant-finite},
\[
 Q:=V/(\varpi,K-1)=L\otimes_{\mathcal E}C_K
\]
is finite over \(D/\varpi\). It suffices, by compact Nakayama over
\(\mathscr A\), to show that \(Q\) is finite-dimensional over \(\kk\).

Suppose otherwise. Its closed support contains a one-dimensional
complete local domain quotient of \(D/\varpi\). At the generic point
of this curve, \(Q\) has nonzero fiber. Since \(D\) is finite over
\(R\), its image in \(\Spec R\) is also a nonconstant curve. Extend
its fraction field to a field \(H\) as in
Lemma~\ref{lem:curve-support}. With subscripts denoting fibers, we have
\[
                 L_H\otimes_{\mathcal E_H}C_{K,H}\ne0.
\]
This tensor product still has the commuting actions of the remaining
\(b-1\) copies of \(\mathcal E_H\).

For any nonzero finite right \(\mathcal E_H\)-module \(U\), choose a
simple quotient \(S\). Right exactness of tensor products and
Lemma~\ref{lem:curve-support} give
\(U\otimes_{\mathcal E_H}C_{K,H}\ne0\). Apply this observation
successively to the remaining factors. The central parameters are
equal, so the same algebra and the same module \(C_{K,H}\) occur each
time. We obtain
\[
 L_H\otimes_{\mathcal E_H^{\otimes_H b}}C_{K,H}^{\otimes_H b}\ne0.
\]
By the product formula in Lemma~\ref{lem:coinvariant-finite}, this is
the fiber of \(M/(\varpi,K^b-1)\) on the chosen curve. But product
admissibility makes that module finite-dimensional over \(\kk\).
Its continuous \(D\)-action is supported at the closed point, so its
fiber on a nonconstant curve is zero. This contradiction proves that
\(Q\) is finite-dimensional and hence that \(V\) is finite over
\(\mathscr A\).
\end{proof}
